\section{Methods and conventions}\label{sec:methods}
We use the physical shift $(Su)(j)=u(j-1)$, so $S^c\delta_0=\delta_c$, and $X\delta_j=j\delta_j$. Reflection is $(\mathcal Ru)(j)=u(-j)$. The potential is $v(t)=2\lambda\cos2\pi t$. All inner products are conjugate-linear in the first argument. Analytic continuations of matrix entries are taken before any conjugation used on the real axis.

Chambers' formula, in the uniform Bloch gauge, is
\begin{equation}\label{eq:Chambers}
 \det(E-B_t(k))=P(E)-2\cos k-c(t),\qquad
 c(t)=2\Lambda\cos2\pi qt.
\end{equation}
Here $B_t(k)$ has cyclic hopping factors $e^{\pm ik/q}$. The sign of $k$ is fixed so that $[A,X](k)=-iq\partial_kA(k)$. The physical shift or its adjoint can be used as the cyclic hopping matrix; this declaration fixes the orientation in the orbital formulas. All energy spectra and all width estimates are invariant under reversing $k$.

Six inputs organize the proof. (I1) is Chambers' identity. (I2) is the near-spectrum derivative estimate in Paper~A, Proposition~3.3(2). (I3) is the complex-energy extension of Paper~B's growth and rational-product lemmas, proved in Appendix~\ref{app:complex}. (I4) is the Aubry--Perron maximal-site fact at $B_0(0)$, proved in Appendix~\ref{app:frames}. (I5) is the common uniform Mathieu ladder, derived from one-well quantization in Appendix~\ref{app:mathieu}. (I6) comprises the continuous IDS, full DOS support, and the convention $N=\{m\alpha\}$; see Avron--Simon~\cite{AvronSimon} and Avila--You--Zhou~\cite[Section~1]{AYZdry}.

For fixed $\xi\in(0,L)$ we set $\tau=\xi/100$. Only finitely many derivatives and resolvent products occur. Let $C_*$ dominate their exponential-loss coefficients. We choose $h>L$ only after that list is fixed, so that
\begin{equation}\label{eq:slack}
 C_*(h-L)<\min\{\xi,\tau,L\}/10000.
\end{equation}
The notation $\EE=e^{o(q)}$ in the appendices denotes a positive quantity bounded by $Cq^Ce^{C(h-L)q}$ for one of this finite list of constants. Each occurrence can be bounded by an arbitrarily small fixed exponential loss by this choice and a subsequent large-$q$ threshold. It never hides dependence on $b,D$, the well or $\theta$. To obtain a literal $o(1)$ in an exponent, diagonalize the fixed-loss choices and their thresholds. The extension of Theorem~\ref{thm:scale} to $\xi=L$, if desired, follows by applying any smaller margin and restricting to its single permitted label $D=1$.

Dirichlet compression to an interval $I$ is denoted by a subscript $I$. Eigenvalues are counted with multiplicity. Strict inequalities in counting windows are used only away from the enclosing clusters. Schur norms always include both row and column sums. Weighted Schur estimates are proved with their spatial weights before the corresponding norm is simplified.

\label{end:methods}
